\begin{center}  \Large{\textbf{ΕΚΤΕΤΑΜΕΝΗ ΕΛΛΗΝΙΚΗ ΠΕΡΙΛΗΨΗ}}  \end{center}
\begin{center}  \Large{\textbf{Επεξεργασία και κωδικοποίηση κίνησης σκελετικών
μοντέλων για την ομαλοποίηση και ανάλυση δεδομένων}}  \end{center}
\vspace{3pt}
\begin{center}  \large{\textbf{Νικόλαος Γέροντας: \hspace{50pt} Κωνσταντίνος Μουστάκας:}}  \end{center}
\vspace{10pt}
\noindent

Η αλληλεπίδραση με τρισδιάστατο περιεχόμενο μέσω των χεριών στηρίζεται
συνήθως σε εξειδικευμένο εξοπλισμό --- κάμερες βάθους, κράνη εικονικής
πραγματικότητας, ειδικά γάντια ή χειριστήρια --- καθένα από τα οποία όμως
αυξάνει το κόστος και τη δυσκολία πρόσβασης στις ίδιες ακριβώς
αλληλεπιδράσεις που υποστηρίζει. Η παρούσα διπλωματική εργασία διερευνά αν
μια αντίστοιχη διαδραστική τρισδιάστατη εφαρμογή, ελεγχόμενη με τα χέρια,
μπορεί να υλοποιηθεί μόνο με μια κάμερα και έναν κοινό
\acrshort{cpu}, χωρίς κράνος, αισθητήρα βάθους ή ειδικά γάντια.

Η δυσκολία έγκειται στην αναπαράσταση. Το Google MediaPipe εντοπίζει 21
σημεία-κλειδιά ανά χέρι σε πραγματικό χρόνο και σε κοινό υλικό, ωστόσο τα
σημεία αυτά αποτελούν μια απλή λίστα συντεταγμένων χωρίς εσωτερική δομή: η
κλίμακά τους εξαρτάται από την απόσταση από την κάμερα, το βάθος αποδίδεται
μόνο ως προς τον καρπό, ενώ δεν περιέχουν καμία πληροφορία για την ιεραρχία
των αρθρώσεων ή για τα ανατομικά επιτρεπτά εύρη κίνησης. Το παραμετρικό
μοντέλο χεριού \acrshort{mano} προσφέρει ακριβώς αυτή τη δομή που λείπει:
ένα σταθερό κινηματικό δέντρο περιστροφών στις αρθρώσεις, ανεξάρτητο από το
μέγεθος του χεριού και την οπτική γωνία. Η σύνδεση των δύο αυτών
αναπαραστάσεων αποτελεί τον πυρήνα της εργασίας.

Τη σύνδεση αυτή υλοποιεί ένας αναλυτικός επιλύτης αντίστροφης κινηματικής,
ο οποίος υπολογίζει τη στάση του χεριού σε κλειστή μορφή: μια ευθυγράμμιση
των συνόλων σημείων μέσω \acrshort{svd} για τον καθολικό προσανατολισμό και,
στη συνέχεια, μία μόνο διάσχιση του κινηματικού δέντρου, κατά την οποία ανακτάται
για κάθε άρθρωση μόνο η συνιστώσα αιώρησης (swing), ενώ η συστροφή
(twist) τίθεται μηδενική. Τα σημεία-κλειδιά εξομαλύνονται
αρχικά με ένα φίλτρο 1\euro{} και, έπειτα, οι υπολογισμένες περιστροφές, η
αναγνωρισμένη χειρονομία και ένα σημείο αναφοράς για την τοποθέτηση
σειριοποιούνται σε ένα πακέτο \acrshort{json} και μεταδίδονται μέσω
\acrshort{udp} σε μια εφαρμογή Unity. Εκεί, ένα κοινό επίπεδο αλληλεπίδρασης
τροφοδοτεί τέσσερις σκηνές, καθεμία από τις οποίες αξιοποιεί την ίδια στάση
με εντελώς διαφορετικό τρόπο: ένα αρκουδάκι που λειτουργεί ως σύντροφος και
ως μέσο αναρρίχησης, ένα πρόσωπο που παραμορφώνεται ανά κορυφή, μια
επεξεργάσιμη πόλη και ένα εργαλείο ζωγραφικής, του οποίου τα σχέδια εξάγονται
από την εφαρμογή ως νέφη σημείων για ανακατασκευή επιφάνειας.

Ολόκληρη η διαδικασία εκτελείται αποκλειστικά στον \acrshort{cpu}, χωρίς να
εμπλέκεται σε κανένα στάδιο η \acrshort{gpu}. Η αξιολόγηση της καθυστέρησης
σε \num{20234} καρέ κατέγραψε μέση καθυστέρηση λογισμικού \SI{30.68}{\ms},
με το 99ο εκατοστημόριο στα \SI{40}{\ms} --- αρκετά χαμηλότερα από το όριο
των \SI{60}{\ms}, πέρα από το οποίο η καθυστέρηση γίνεται αντιληπτή από τον
χρήστη. Το σύστημα θυσιάζει έτσι συνειδητά την ακρίβεια της παρακολούθησης
χάριν της προσβασιμότητας, προσφέροντας τρισδιάστατες αλληλεπιδράσεις
ακριβείας με τα δύο χέρια, με διαδραστική απόκριση, μέσα από την ενσωματωμένη
κάμερα ενός κοινού φορητού υπολογιστή.\\\\



Ακολουθεί περίληψή κάθε κεφαλαίου που αναπτύχθηκε, στα Ελληνικά:

\subsection*{Κεφάλαιο 1: Εισαγωγή}

Η εργασία ξεκινά από ένα απλό ερώτημα: είναι δυνατόν να δημιουργηθεί μια
εφαρμογή που συνδυάζει το παιχνίδι με την επεξεργασία περιεχομένου και
ελέγχεται από τις κινήσεις των χεριών, χωρίς κανέναν επιπλέον εξοπλισμό ---
χωρίς κράνος εικονικής πραγματικότητας, κάμερα βάθους ή ειδικά γάντια ---
παρά μόνο με την κάμερα που διαθέτει ήδη ένας φορητός υπολογιστής. Το κίνητρο
είναι η προσβασιμότητα: τα ακριβά περιφερειακά περιορίζουν το ποιος μπορεί να
χρησιμοποιήσει ένα σύστημα, ενώ το χέρι αποτελεί τον πιο φυσικό τρόπο ελέγχου.

Η προσέγγιση αυτή έγινε εφικτή χάρη σε εργαλεία όπως το MediaPipe, το οποίο
εντοπίζει σε πραγματικό χρόνο 21 σημεία-κλειδιά ανά χέρι από μια απλή έγχρωμη
κάμερα. Τα ακατέργαστα αυτά σημεία, ωστόσο, δεν αρκούν για να κινήσουν έναν
σκελετό ή να καταγράψουν την κίνηση με νόημα: η κλίμακά τους αλλάζει ανάλογα
με την απόσταση από την κάμερα, το βάθος εκτιμάται αντί να μετριέται και δεν
περιέχουν καμία πληροφορία για την ιεραρχία των αρθρώσεων. Απαιτείται μια
κανονικοποιημένη αναπαράσταση, την οποία προσφέρει το παραμετρικό μοντέλο
\acrshort{mano}, περιγράφοντας το χέρι ως ένα σύνολο περιστροφών αρθρώσεων
πάνω σε ένα σταθερό κινηματικό δέντρο.

Το κεντρικό πρόβλημα της εργασίας είναι, επομένως, η σύνδεση της εξόδου του
MediaPipe με την αναπαράσταση του \acrshort{mano}, σε πραγματικό χρόνο, σε
επεξεργαστή και με αρκετά χαμηλή καθυστέρηση ώστε η εφαρμογή να παραμένει
αποκρίσιμη. Η συνεισφορά εκτείνεται σε ολόκληρη τη διαδρομή, από την εικόνα
της κάμερας έως μια εφαρμογή Unity: ένα υποσύστημα παρακολούθησης και εξαγωγής
τρισδιάστατων συντεταγμένων με εξομάλυνση μέσω του φίλτρου 1\euro{}, έναν
αναλυτικό επιλύτη αντίστροφης κινηματικής που αντιστοιχίζει τα 21 σημεία στον
σκελετό του \acrshort{mano}, ένα επίπεδο μετάδοσης μέσω \acrshort{udp}, και
ένα πρωτότυπο σε Unity με τέσσερις διαφορετικές σκηνές.



\subsection*{Κεφάλαιο 2: Θεωρητικό Υπόβαθρο}

Το δεύτερο κεφάλαιο θεμελιώνει θεωρητικά τα επιμέρους τμήματα του συστήματος,
με τη σειρά που αυτά εμφανίζονται στη ροή επεξεργασίας. Ξεκινά από την
ανατομία του χεριού --- τα 27 οστά, τις αρθρώσεις και τους βαθμούς ελευθερίας
τους --- και την αναπαράστασή του ως κινηματικής αλυσίδας, δηλαδή ως δέντρου
άκαμπτων τμημάτων που συνδέονται με αρθρώσεις. Παρουσιάζονται επίσης οι δύο
διαφορετικές συμβάσεις αρίθμησης που χρησιμοποιούνται στη συνέχεια: τα 21
σημεία-κλειδιά του MediaPipe και οι 16 αρθρώσεις του \acrshort{mano}.

Ακολουθεί η αναλυτική περιγραφή του μοντέλου \acrshort{mano}, ενός στατιστικού
μοντέλου του ανθρώπινου χεριού που έχει εκπαιδευτεί από τρισδιάστατες σαρώσεις
και βασίζεται στη λογική του μοντέλου σώματος SMPL. Το μοντέλο διαχωρίζει τις
μεταβολές της γεωμετρίας σε εκείνες που οφείλονται στην ταυτότητα του ατόμου
και σε εκείνες που προκαλούνται από τη στάση, ενώ οι διορθωτικές συνιστώσες
στάσης αντιμετωπίζουν την κατάρρευση του πλέγματος στις αρθρώσεις, την οποία
εμφανίζει το απλό γραμμικό skinning. Στη συνέχεια περιγράφεται το MediaPipe ως
μια διεργασία δύο σταδίων --- ανίχνευση της παλάμης και εντοπισμός των σημείων
--- και ο τρόπος με τον οποίο αναγνωρίζει χειρονομίες από την κατάσταση των
δαχτύλων.

Το κεφάλαιο αναπτύσσει έπειτα τον πυρήνα της μεθόδου, την αντίστροφη
κινηματική. Διακρίνει την ευθεία από την αντίστροφη κινηματική και εξηγεί
γιατί, όταν όλες οι θέσεις των αρθρώσεων είναι ταυτόχρονα γνωστές, το πρόβλημα
λύνεται σε κλειστή μορφή, χωρίς επαναλήψεις. Ο αλγόριθμος AIK ανακτά τον
καθολικό προσανατολισμό του καρπού μέσω μίας ευθυγράμμισης συνόλων σημείων με
ανάλυση ιδιαζουσών τιμών (\acrshort{svd}) και, στη συνέχεια, υπολογίζει κάθε
επιμέρους περιστροφή αναλυτικά, με μια αποσύνθεση αιώρησης--συστροφής
(swing--twist) κατά μήκος της αλυσίδας, θέτοντας τη συστροφή ίση με το μηδέν
--- μια απλοποίηση που είναι φυσιολογικά αμελητέα και ταυτόχρονα αίρει την
ασάφεια κατεύθυνσης όταν τα οστά είναι σχεδόν συγγραμμικά.

Τέλος, καλύπτονται τα επικοινωνιακά και υποστηρικτικά τμήματα: η επιλογή του
πρωτοκόλλου \acrshort{udp} έναντι του TCP για τη μετάδοση σε πραγματικό χρόνο
--- το \acrshort{udp} δεν περιμένει ποτέ να ανακτήσει χαμένα πακέτα, οπότε το
χειρότερο κόστος μιας απώλειας είναι ένα μόνο χαμένο καρέ αντί για αύξηση της
καθυστέρησης --- και η σειριοποίηση των δεδομένων σε μορφή \acrshort{json}· το
φίλτρο 1\euro{}, το οποίο εξισορροπεί την τρεμούλα (jitter) με την καθυστέρηση
(lag) προσαρμόζοντας τη συχνότητα αποκοπής στην ταχύτητα της κίνησης· η μηχανή
Unity· και η ανακατασκευή επιφανειών από νέφη σημείων με τη νευρωνική μέθοδο NKSR.



\subsection*{Κεφάλαιο 3: Σύγχρονη Τεχνολογική Στάθμη}

Το τρίτο κεφάλαιο τοποθετεί την εργασία μέσα στο ευρύτερο τοπίο των
ανταγωνιστικών προσεγγίσεων, κρίνοντάς τες με βάση τους τρεις περιορισμούς που
τη διέπουν: μονοφθαλμική είσοδο RGB, εκτέλεση σε πραγματικό χρόνο μόνο σε
επεξεργαστή και χρήση κοινού, εμπορικά διαθέσιμου εξοπλισμού. Μεγάλο μέρος των
ακριβέστερων μεθόδων προϋποθέτει κάμερα βάθους, υπολογιστική ισχύ επιπέδου
κάρτας γραφικών ή διατάξεις πολλαπλών καμερών στην κεφαλή· τέτοιες μέθοδοι δεν
απορρίπτονται, αλλά σταθμίζονται έναντι του κόστους που θα επέβαλλαν στον μέσο
χρήστη.

Ως προς την εκτίμηση της στάσης και του σχήματος, οι μέθοδοι χωρίζονται σε δύο
οικογένειες. Οι μέθοδοι χωρίς μοντέλο εκπαιδεύουν ένα νευρωνικό δίκτυο να
προβλέπει απευθείας τις συντεταγμένες των αρθρώσεων· είναι ταχείες, αλλά η
έξοδός τους είναι ένα άμορφο σύνολο σημείων, χωρίς περιστροφές ή ιεραρχία. Οι
μέθοδοι με μοντέλο, αντίθετα, στηρίζουν τις προβλέψεις τους στο \acrshort{mano}
και παράγουν ανατομικά έγκυρο αποτέλεσμα· χαρακτηριστικά παραδείγματα είναι τα
συστήματα των Zhou κ.ά., το HandOccNet για ανθεκτικότητα στις αποκρύψεις και
το HaMeR. Τα συστήματα αυτά, ωστόσο, αξιολογούνται σχεδόν αποκλειστικά σε
κάρτες γραφικών. Ακόμη και μέθοδοι που τρέχουν σε πραγματικό χρόνο σε φορητές
συσκευές, όπως το MobRecon, προβλέπουν απευθείας τις κορυφές του πλέγματος και
όχι περιστροφές, οπότε δεν μπορούν να οδηγήσουν έναν σκελετό χωρίς ένα πρόσθετο
βήμα.

Στη συνέχεια συγκρίνονται τα παραμετρικά μοντέλα χεριού. Το \acrshort{mano}
αποτελεί το de facto πρότυπο του πεδίου, χάρη στον ελαφρύ και κατάλληλο για
πραγματικό χρόνο υπολογισμό του, ενώ το NIMBLE, παρότι μοντελοποιεί την
εσωτερική ανατομία --- οστά, μυς και δέρμα --- με στόχο τη φωτορεαλιστική
απόδοση, απαιτεί κάρτα γραφικών και προορίζεται για επεξεργασία εκτός
πραγματικού χρόνου. Για τους λόγους αυτούς υιοθετείται το \acrshort{mano}.

Το κεφάλαιο εξετάζει έπειτα τις τέσσερις οικογένειες επιλυτών αντίστροφης
κινηματικής. Οι επαναληπτικοί αριθμητικοί επιλύτες φέρουν περιττό κόστος
σύγκλισης, αφού η είσοδος είναι πλήρως παρατηρούμενη· η προσαρμογή μοντέλου με
βελτιστοποίηση παράγει ανατομικά ομαλά αποτελέσματα, αλλά με μη φραγμένο χρόνο
επίλυσης, ακατάλληλο για διαδραστική χρήση σε επεξεργαστή· η νευρωνική
παλινδρόμηση, και ιδίως το HybrIK-X, πετυχαίνει τη μεγαλύτερη ακρίβεια
ανακτώντας τη συνιστώσα της συστροφής που η γεωμετρία από μόνη της δεν μπορεί
να προσδιορίσει, αλλά με κόστος μια κάρτα γραφικών. Ο αναλυτικός επιλύτης AIK
είναι ο μόνος που ικανοποιεί ταυτόχρονα και τους τρεις περιορισμούς: δεν
εξαρτάται από δεδομένα, εκτελείται σε ντετερμινιστικό χρόνο σε επεξεργαστή και
έχει εξ ορισμού φραγμένο κόστος ανά καρέ. Έτσι, το HybrIK-X ορίζει το ανώτατο
όριο ακρίβειας που το σύστημα δεν φτάνει, ενώ το AIK ορίζει το κατώτατο όριο
καθυστέρησης και υλικού εντός του οποίου αυτό λειτουργεί.

Τέλος, εξετάζονται η αλληλεπίδραση και οι δημιουργικές εφαρμογές χωρίς δείκτες
(markers). Από το ιστορικό \emph{Put-that-there} του Bolt και τις πρώτες
οπτικές προσεγγίσεις, έως τα παιχνίδια και τα εμβυθιστικά περιβάλλοντα (Leap
Motion, MEgATrack στο Oculus Quest) και τα εργαλεία τρισδιάστατης δημιουργίας
(CavePainting, Tilt Brush, Gravity Sketch), το κοινό μοτίβο είναι ότι τα πιο
ικανά συστήματα στηρίζονται σε ειδικό εξοπλισμό --- οπτικά περιφερειακά, κράνη
ή χειριστήρια. Το εργαλείο ζωγραφικής της παρούσας εργασίας προσφέρει την ίδια
ακριβώς δυνατότητα --- την εναπόθεση γεωμετρίας στον χώρο με την κίνηση του
χεριού --- χωρίς όμως τον εξοπλισμό που οι μέθοδοι αυτές προϋποθέτουν.



\subsection*{Κεφάλαιο 4: Σύστημα Παρακολούθησης Χεριών σε Πραγματικό Χρόνο}

Το τέταρτο κεφάλαιο περιγράφει την υλοποίηση του συστήματος σε Python, δηλαδή
τη διαδικασία που μετατρέπει ένα μεμονωμένο καρέ της κάμερας σε μια
κανονικοποιημένη στάση \acrshort{mano} και τη μεταδίδει, μαζί με την
αναγνωρισμένη χειρονομία και ένα σημείο τοποθέτησης, σε μια ξεχωριστή διεργασία
απόδοσης.

Η υλοποίηση χωρίζεται σε δύο στάδια. Ένα στάδιο εκτός σύνδεσης εκτελείται μία
φορά και μετατρέπει την επίσημη διανομή του \acrshort{mano} σε μια ελαφριά
μορφή, που φορτώνεται γρήγορα και δεν εξαρτάται από παρωχημένες βιβλιοθήκες. Το
διαδικτυακό στάδιο επεξεργάζεται κάθε καρέ σε πραγματικό χρόνο και περιλαμβάνει
ολόκληρη τη διαδρομή των δεδομένων, από τη σύλληψη έως τη μετάδοση. Ο
κινηματικός πυρήνας συγκεντρώνεται σε ένα πακέτο, του οποίου τα επιμέρους
τμήματα αντιστοιχούν άμεσα στις θεωρητικές έννοιες του δεύτερου κεφαλαίου, ενώ
γύρω από αυτό ένα σύνολο μονάδων ανώτερου επιπέδου οδηγεί τη ζωντανή ροή.

Η ανίχνευση γίνεται με τον αναγνωριστή χειρονομιών του MediaPipe, ρυθμισμένο να
παρακολουθεί έως δύο χέρια. Λειτουργεί ασύγχρονα: ο βρόχος σύλληψης στέλνει
κάθε καρέ και διαβάζει αμέσως όποιο αποτέλεσμα είναι διαθέσιμο, χωρίς ποτέ να
μπλοκάρει περιμένοντας τον ανιχνευτή. Για κάθε χέρι, το MediaPipe επιστρέφει
δύο σύνολα σημείων: τα σημεία στον τρισδιάστατο χώρο, εκφρασμένα σε μετρικές
μονάδες με αρχή τον καρπό, που περιγράφουν το σχήμα του χεριού, και τα σημεία
στο επίπεδο της εικόνας, που καθορίζουν τη θέση του. Τα πρώτα εξομαλύνονται με
το φίλτρο 1\euro{}, ενώ από τα δεύτερα προκύπτει ένα σημείο τοποθέτησης, που
συνδυάζει τη θέση του καρπού στην εικόνα με μια εκτίμηση του βάθους από τον
λόγο των μηκών των οστών.

Ο επιλύτης δέχεται έναν σκελετό προτύπου και έναν σκελετό-στόχο και επιστρέφει
τις περιστροφές που τοποθετούν το πρότυπο πάνω στον στόχο. Ανακτά πρώτα τον
καθολικό προσανατολισμό του καρπού μέσω \acrshort{svd} και έπειτα διασχίζει το
κινηματικό δέντρο, υπολογίζοντας για κάθε άρθρωση την περιστροφή αιώρησης. Όλη
η μετάβαση από τα σημεία στις περιστροφές ανάγεται σε μία ανάλυση πίνακα,
ακολουθούμενη από μία μόνο διάσχιση του δέντρου, χωρίς καμία επανάληψη και
χωρίς εκπαιδευμένο δίκτυο στη διαδρομή. Οι περιστροφές μετατρέπονται σε
διανύσματα άξονα-γωνίας, κανονικοποιούνται ως προς την κλίμακα του προτύπου και
εφαρμόζονται στο μοντέλο για την παραμόρφωση του πλέγματος.

Τέλος, η στάση σειριοποιείται σε ένα πακέτο \acrshort{json}, που περιέχει τη
χρονική πληροφορία και, για κάθε χέρι, τις δεκαέξι περιστροφές, τη χειρονομία και το
σημείο τοποθέτησης· τα πεδία ενός χεριού που δεν εντοπίστηκε παραμένουν κενά.
Το πακέτο αποστέλλεται μέσω \acrshort{udp} στην τοπική διεύθυνση, και μόνο
εφόσον έχει εντοπιστεί τουλάχιστον ένα χέρι, ώστε το κανάλι να παραμένει
ανενεργό όταν δεν υπάρχει ορατό χέρι.



\subsection*{Κεφάλαιο 5: Η Διαδραστική Εφαρμογή Unity}

Το πέμπτο κεφάλαιο περιγράφει την εφαρμογή Unity που χρησιμοποιεί τη ροή.
Ολόκληρη η εφαρμογή οργανώνεται γύρω από μία αρχή: ο δέκτης είναι το μόνο τμήμα
που επικοινωνεί με το δίκτυο, ενώ κάθε σκηνή απλώς διαβάζει σε κάθε καρέ τον
τρέχοντα σκελετό, τη χειρονομία και το σημείο τοποθέτησης. Ο ανακατασκευασμένος
σκελετός λειτουργεί έτσι ως κοινός πόρος, και οι σκηνές παραμένουν ανεξάρτητες
από το υποκείμενο πρωτόκολλο.

Ο δέκτης αντιστρέφει βήμα προς βήμα τη σειριοποίηση: τρέχει σε ξεχωριστό νήμα,
αποκωδικοποιεί κάθε πακέτο και μετατρέπει τα διανύσματα άξονα-γωνίας σε στροφές,
που εφαρμόζονται στις δεκαέξι αρθρώσεις του σκελετού, ενώ το σημείο τοποθέτησης
ορίζει τη θέση ολόκληρου του χεριού στη σκηνή. Επειδή σε κάθε καρέ μετριέται η
διαφορά ανάμεσα στη στιγμή αποστολής και στη στιγμή λήψης, ο δέκτης είναι και
το σημείο όπου καταγράφεται η καθυστέρηση του λογισμικού.

Πάνω σε αυτή τη βάση χτίζεται ένα κοινό επίπεδο αλληλεπίδρασης. Η εφαρμογή
χρησιμοποιεί πέντε διακριτές χειρονομίες, με τις οποίες ο χρήστης πλοηγείται
ανάμεσα στις σκηνές, χειρίζεται ένα μενού προσαρτημένο στον καρπό του και
σημαδεύει αντικείμενα στον χώρο, χωρίς ποντίκι ή κουμπιά. Ένας μικρός
μηχανισμός ανοχής απορροφά τις στιγμιαίες απώλειες μιας χειρονομίας από τον
ανιχνευτή, ώστε ένα συνεχές κράτημα να μην τρεμοπαίζει.

Ακολουθούν οι τέσσερις σκηνές. Στην πρώτη, ένα αρκουδάκι βρίσκεται κάτω από δύο
καθεστώτα ελέγχου: άλλοτε λειτουργεί ως σύντροφος που αντιδρά στον χρήστη και
άλλοτε ως avatar που ελέγχει ο ίδιος ο χρήστης --- με κλειστή γροθιά εκτοξεύει
σχοινιά, αιωρείται μέσα στη σκηνή και κατευθύνει την κάμερα με το ελεύθερο
χέρι. Η δεύτερη σκηνή, μια αναδημιουργία του διαδραστικού προσώπου από μια
παλιά κονσόλα, αποτελεί τον πιο άμεσο έλεγχο της χωρικής ακρίβειας: ο χρήστης
πιάνει και παραμορφώνει ένα πλέγμα κεφαλιού κορυφή προς κορυφή, με μια ομαλή
κατανομή βάρους γύρω από το σημείο της λαβής και μια ελαστική επαναφορά τύπου
μάζας--ελατηρίου. Η τρίτη σκηνή είναι μια επεξεργάσιμη πόλη, όπου ο χρήστης
επιλέγει μέσα από το μενού του καρπού ανάμεσα στην πλοήγηση, τη μετακίνηση των
κτιρίων με τα χέρια και την εναέρια μετακίνηση με διπλά σχοινιά. Η τέταρτη
σκηνή είναι το δημιουργικό εργαλείο: ο χρήστης ζωγραφίζει εκτοξεύοντας μια
ακτίνα από την άκρη του δείκτη πάνω σε μια επιφάνεια, μπορεί να σχεδιάσει είτε
σε δύο διαστάσεις είτε --- αλλάζοντας οπτική γωνία --- να συσσωρεύσει μια
τρισδιάστατη μορφή ή να ορίσει ένα στερεό εκ περιστροφής (lathe), και τελικά να
εξάγει το σχέδιο ως νέφος σημείων, που τροφοδοτεί το στάδιο της ανακατασκευής
επιφανειών.

Το κεφάλαιο κλείνει συνοψίζοντας την κοινή παρατήρηση που διατρέχει και
τις τέσσερις σκηνές: κάθε αλληλεπίδραση ανήκει σε ένα είδος που τα
αντίστοιχα συστήματα προσφέρουν μόνο μέσω ειδικού περιφερειακού, κράνους ή
χειριστηρίου, ενώ εδώ επιτυγχάνεται από μία και μόνη κάμερα και έναν κοινό
\acrshort{cpu}, με τον χρήστη να μην κρατά και να μη φορά τίποτα πέρα,
προαιρετικά, από ένα απτικό γιλέκο. Η αξιολόγηση που ακολουθεί στο επόμενο
κεφάλαιο επιβεβαιώνει ότι ο βρόχος δεν είναι απλώς κλειστός, αλλά αρκετά
γρήγορος και οικονομικός ώστε οι αλληλεπιδράσεις αυτές να γίνονται αισθητές
ως άμεσες.



\subsection*{Κεφάλαιο 6: Το Τελικό Σύστημα και η Αξιολόγησή του}

Το έκτο κεφάλαιο συνθέτει πρώτα την εργασία σε ένα ενιαίο σύνολο και έπειτα
τη θέτει υπό δοκιμή. Οι δύο διεργασίες που περιγράφηκαν χωριστά στα προηγούμενα
κεφάλαια — η ροή σε Python που μετατρέπει κάθε καρέ σε ένα μεταδιδόμενο πακέτο
και η εφαρμογή Unity που μετατρέπει τη ροή αυτή σε αλληλεπίδραση — εκτελούμενες
μαζί συνιστούν ένα και μόνο σύστημα: μια εφαρμογή παρακολούθησης χεριών και
αλληλεπίδρασης σε πραγματικό χρόνο, με μία απλή κάμερα και αποκλειστικά σε \acrshort{cpu},
μέσα στην οποία ο χρήστης τοποθετεί τα γυμνά του χέρια σε μια τρισδιάστατη σκηνή.
Αυτό που το καθιστά σύστημα και όχι απλή επίδειξη είναι ότι ο ίδιος βρόχος
τροφοδοτεί, χωρίς καμία τροποποίηση, και τις τέσσερις διαφορετικές σκηνές. Πάνω σε αυτό
το ενιαίο σύστημα, το κεφάλαιο εξετάζει ποσοτικά τον κεντρικό ισχυρισμό της εργασίας: ότι
η παρακολούθηση χεριών μόνο με κάμερα και αποκλειστικά σε \acrshort{cpu}
μπορεί να οδηγήσει αλληλεπίδραση που γίνεται αισθητή ως άμεση. Για να σταθεί
ο ισχυρισμός αυτός, πρέπει να ισχύουν δύο ιδιότητες: η ανακτώμενη στάση να
φτάνει στην εφαρμογή αρκετά γρήγορα ώστε η κίνηση να μη γίνεται αντιληπτή ως
καθυστερημένη, και το κόστος παραγωγής της να είναι αρκετά μικρό ώστε
ολόκληρο το σύστημα να εκτελείται σε κοινό υλικό, χωρίς προσφυγή σε κάρτα
γραφικών. Για τον σκοπό αυτό χρησιμοποιούνται τρία εργαλεία μέτρησης: η
ενσωματωμένη στον δέκτη καταγραφή της καθυστέρησης ανά καρέ σε μια μακρά
συνεδρία· μια ξεχωριστή, πιο λεπτομερής διαδρομή χρονομέτρησης που απομονώνει
τα τρία υπολογιστικά στάδια της διεργασίας αποστολής· και ένας παρακολουθητής
πόρων σε επίπεδο διεργασίας, που δειγματοληπτεί τον επεξεργαστή και τη μνήμη
και των δύο διεργασιών.

Ως προς την καθυστέρηση του λογισμικού, σε μια συνεχή συνεδρία \num{20234}
καρέ το σύστημα κατέγραψε μέση καθυστέρηση \SI{30.68}{\ms}, με διάμεσο
\SI{31}{\ms}, τυπική απόκλιση \SI{5.03}{\ms} και το 95ο και 99ο
εκατοστημόριο στα \SI{37}{\ms} και \SI{40}{\ms} αντίστοιχα. Σχεδόν όλα τα
καρέ βρίσκονται έτσι αρκετά κάτω από το όριο των \SI{60}{\ms}, πάνω από το
οποίο η καθυστέρηση γίνεται αντιληπτή· η μέγιστη τιμή των \SI{428}{\ms} ήταν
ένα μεμονωμένο ακραίο συμβάν, συμβατό με μια διακοπή χρονοπρογραμματισμού του
λειτουργικού ή με μια παύση συλλογής σκουπιδιών, χωρίς καμία διαρκή
υποβάθμιση κατά τη διάρκεια της συνεδρίας.

Η ανάλυση της καθυστέρησης ανά στάδιο, σε \num{5320} καρέ, δείχνει πού
δαπανάται ο χρόνος: ο αναλυτικός επιλύτης κυριαρχεί με το \SI{89.7}{\percent}
του υπολογισμού ανά καρέ, η ανίχνευση του MediaPipe αντιστοιχεί στο
\SI{8.1}{\percent} και το στάδιο εξομάλυνσης και αγκύρωσης μόλις στο
\SI{1.9}{\percent}, ενώ το υπόλοιπο \SI{0.3}{\percent} αποδίδεται στη
διαχείριση ανά χέρι γύρω από τον επιλύτη. Πρακτικά, ο επιλύτης είναι το μόνο
τμήμα του οποίου το κόστος έχει σημασία· το γεγονός αυτό δικαιώνει την επιλογή
να διατηρηθούν σκόπιμα ελαφριά η ανίχνευση και το φίλτρο 1\euro{} και να
συγκεντρωθεί η αναλυτική εργασία στο στάδιο ανάκτησης της στάσης, και
υποδεικνύει ότι κάθε μελλοντική προσπάθεια μείωσης της καθυστέρησης --- ή
υποστήριξης περισσότερων ταυτόχρονων χεριών --- θα έπρεπε να εστιάσει εκεί
και σχεδόν πουθενά αλλού.

Τέλος, ως προς τους πόρους, η διεργασία αποστολής --- που φέρει ολόκληρη τη
ροή όρασης και κινηματικής --- κατανάλωσε κατά μέσο όρο το \SI{17.5}{\percent}
της συνολικής χωρητικότητας του επεξεργαστή και η εφαρμογή Unity το
\SI{14.6}{\percent}· μαζί καταλαμβάνουν περίπου το ένα τρίτο της μηχανής. Η
μνήμη παρέμεινε χαμηλή και, κυρίως, σταθερή --- περίπου \num{495}\,MB και
\num{481}\,MB αντίστοιχα, χωρίς ουσιαστική αύξηση κατά τη διάρκεια της
συνεδρίας --- η αναμενόμενη συνέπεια μιας σχεδίασης όπου διατηρείται πάντοτε
μόνο το πιο πρόσφατο πακέτο και δεν συσσωρεύεται καμία κατάσταση ανά καρέ. Το
πλήρες σύστημα των δύο διεργασιών εκτελείται έτσι άνετα εντός του
προϋπολογισμού ενός κοινού φορητού υπολογιστή, χωρίς κάρτα γραφικών και
αφήνοντας, κατά μέσο όρο, πάνω από το \SI{40}{\percent} του επεξεργαστή
αδρανές. Οι τρεις μετρήσεις συγκλίνουν στο ίδιο συμπέρασμα: ο βρόχος δεν είναι
απλώς κλειστός, αλλά γρήγορος και οικονομικός --- και αυτός ακριβώς ο
συνδυασμός, και όχι καθεμία ιδιότητα χωριστά, καθιστά τον ισχυρισμό
προσβασιμότητας της εργασίας μετρήσιμο αποτέλεσμα και όχι απλή φιλοδοξία.



\subsection*{Κεφάλαιο 7: Συμπεράσματα και Μελλοντικές Επεκτάσεις}

Το έβδομο κεφάλαιο απαντά καταφατικά στο αρχικό ερώτημα: μια διαδραστική
τρισδιάστατη εφαρμογή ελεγχόμενη με τα χέρια μπορεί πράγματι να υλοποιηθεί μόνο
με μια κάμερα και έναν επεξεργαστή. Συνοψίζει τη συνεισφορά --- τη σύνδεση του
MediaPipe με το \acrshort{mano} μέσω του AIK, τη ροή σε Python, την εφαρμογή
Unity με τις τέσσερις σκηνές και το ποσοτικό αποτέλεσμα της καθυστέρησης ---
και, στη συνέχεια, εκθέτει τους περιορισμούς.

Ο σημαντικότερος περιορισμός είναι ότι η ανακτώμενη στάση αξιολογήθηκε ως προς
την καθυστέρηση, αλλά όχι ως προς την ακρίβεια: κανένα μετρικό μέγεθος δεν
ποσοτικοποιεί πόσο πιστά αναπαράγει τις πραγματικές περιστροφές του χεριού.
Δεύτερον, ο μηδενισμός της συστροφής, αν και φυσιολογικά θεμελιωμένος, χάνει
ακριβώς τον βαθμό ελευθερίας που έχει σημασία σε ορισμένες στάσεις, όπως η
αντίθεση του αντίχειρα. Επιπλέον, το βάθος δεν μετριέται αλλά εκτιμάται, και
μια μόνο κάμερα δεν προσφέρει καμία λύση όταν το χέρι βγαίνει από το πλαίσιο ή
στρέφεται μακριά της, ενώ η αλληλεπίδραση περιορίζεται στις πέντε χειρονομίες
και στα δύο χέρια.

Από κάθε περιορισμό προκύπτει και μια αντίστοιχη κατεύθυνση για μελλοντική
εργασία: η ποσοτική αξιολόγηση της ακρίβειας σε καθιερωμένα σύνολα δεδομένων· η
προσθήκη ενός ελαφριού εκπαιδευμένου τμήματος που θα ανακτά τη συνιστώσα της
συστροφής, εκτελούμενο μόνο όταν υπάρχει διαθέσιμη κάρτα γραφικών· η ενσωμάτωση
μιας δεύτερης κάμερας ή ενός αισθητήρα βάθους για την αποκατάσταση της μετρικής
κλίμακας· η εκτίμηση του σχήματος ανά χρήστη και η μεταφορά του πλήρους
πλέγματος του \acrshort{mano} μέχρι το στάδιο της απόδοσης· ένα αρχικό στάδιο
επεξεργασίας πιο ανθεκτικό στις αποκρύψεις· ένα προσαρμόσιμο λεξιλόγιο
χειρονομιών· καθώς και η ενσωμάτωση της ανακατασκευής επιφανειών μέσα στην ίδια
την εφαρμογή και η επέκταση της μετάδοσης σε δίκτυο, για απομακρυσμένη ή
πολυχρηστική αλληλεπίδραση.

Το νήμα που διατρέχει τόσο τον σχεδιασμό όσο και τους περιορισμούς είναι μία
και η αυτή συνειδητή επιλογή: η ακρίβεια και η πιστότητα της παρακολούθησης
θυσιάζονται χάριν της προσβασιμότητας. Τα πιο ικανά συστήματα κατακτούν την
ακρίβειά τους προσθέτοντας υλικό και, έτσι, ανεβάζουν το κόστος εισόδου στις
ίδιες τις αλληλεπιδράσεις που προσφέρουν· η παρούσα εργασία ακολουθεί την
αντίθετη λογική, αποδεχόμενη το όριο που επιβάλλει μια απλή κάμερα και μη
ζητώντας από τον χρήστη τίποτα που δεν διαθέτει ήδη. Αυτό που δείχνουν συνολικά
οι τέσσερις σκηνές είναι ότι το όριο αυτό είναι αρκετά υψηλό ώστε να είναι
χρήσιμο: αλληλεπιδράσεις ακριβείας, τρισδιάστατες και με τα δύο χέρια, σε
διαδραστικούς χρόνους, από την κάμερα ενός κοινού φορητού υπολογιστή.
